\section{Where the inefficiency goes: loss categories}\label{app:losses}
\sectionstate{partial}{\texttt{rerun/effbudget/pass1} (\texttt{a1\_recat.py} classes, \texttt{a1\_table.py} table) on the two-track re-pass tables of the chamber~A golden runs (2026-11-03); recovery tests quoted from the pre-pass satstudy are marked}
\intent{Every reference track crossing the chamber~A interior that is not detected, followed to the one place it ends up: one real
event display per class, the fraction of tracks lost to it, its mechanism, and whether a reference-free change to the
reconstruction recovers it (tested by re-running the reconstruction, not argued). Dylan 2026-10-26: a full appendix.}

\begin{table}[tbp]
  \centering\small
  \caption{chamber~A interior, 26\,542 reference tracks of Runs 1 and 2 (one clean reference ray through the interior; Run~1: 6\,454),
  92.29\,\% detected (Run~1: 92.64\,\%): both views fitted, $|r|\le5$\,mm from the reference, $\le2$ competing strips per
  view. Two-track re-pass tables (\texttt{Analysis\_pass1\_twotrack}), classes of \texttt{rerun/effbudget/pass1/a1\_recat.py},
  table generated by \texttt{rerun/effbudget/pass1/a1\_table.py}. Every lost track is in exactly
  one class; first matching class wins, in the order of the table. ``Multi-particle'' = chamber~A alone sees $\ge2$
  particle-like clusters (signal $\le$ one drift column) in a view. The displays below are Run~1 events.}
  \label{tab:losses}
  \begin{tabular}{p{5.3cm}rrp{6.6cm}}
    \hline
    class (figure) & \shortstack{Runs 1+2\\{}[\%]} & \shortstack{Run 1\\{}[\%]} & recoverable without the reference? \\
    \hline
% generated by rerun/effbudget/pass1/a1_table.py -- do not edit
    noise band on a strip patch (\ref{fig:loss_band}) & 0.22 & 0.22 & largely: band hits are censored in the pass (Run 1: 1.43\,\% pre-pass, 0.22\,\% now) \\
    second particle, reference particle fitted, view cut (\ref{fig:loss_multi_m3picked}) & 1.15 & 1.02 & yes, if a view with two tracks is kept (multi-track credit) \\
    second particle, the other one fitted (\ref{fig:loss_multi_otherpicked}) & 1.51 & 1.47 & no: needs a rule to choose the triggering track \\
    second particle, merged in projection (\ref{fig:loss_multi_merged}) & 0.46 & 0.53 & partly: joint two-track fit with charge corroboration, 25/146 of the pass-1 events of this class have a reference-matched fitted track (Run 1: 3/37) \\
    ``spark'': discharge (\ref{fig:loss_spark_discharge}) & 0.50 & 0.50 & no \\
    ``spark'': muon + unrelated deposit (\ref{fig:loss_spark_muonother}) & 0.35 & 0.40 & partly: veto off + competing-cluster cut relaxed \\
    ``spark'': shower at the crossing (\ref{fig:loss_spark_shower}) & 0.81 & 0.67 & no \\
    $\delta$-ray on the track (\ref{fig:loss_delta}) & 0.87 & 0.90 & not with a straight-line model \\
    saturated $\delta$ blob, fit pulled off (\ref{fig:loss_satblob}) & 0.93 & 0.84 & no: neither censoring nor a track\,+\,compact-deposit fit (test on 350 saturated views of Run 1 on the re-pass: 71\,$\to$\,73\,\% agree) \\
    fit off, no other flag (\ref{fig:loss_fit}) & 0.28 & 0.23 & no \\
    fit failed in one view (the pre-pass head-on class) (\ref{fig:loss_headon}) & 0.02 & 0.02 & yes: head-on seeding in the pass (Run 1: 0.26\,\% pre-pass, 17 tracks; 0.02\,\% now, none head-on) \\
    position tail, mostly bent reference tracks (\ref{fig:loss_postail}) & 0.17 & 0.20 & no (reference) \\
    edge of the other coordinate (\ref{fig:loss_edge}) & 0.20 & 0.15 & no (geometry) \\
    second deposit, not particle-like (\ref{fig:loss_multi_deposit}) & 0.21 & 0.19 & no \\
    other (cuts, charge elsewhere)  & 0.03 & 0.03 & --- \\
    \hline
% ends with the closing \hline
  \end{tabular}
\end{table}

\tbw{One paragraph per class: mechanism, how it was identified, the recovery test.}

Each display (\crefrange{fig:loss_band}{fig:loss_multi_deposit}) shows one real chamber~A event of the class in both
views: strip position against drift time, the amplitude of every sample above 4\,\% of the view's maximum, the reference
track extrapolated through the drift gap (dashed) and the production fit (solid). The event is the one used in the
loss-category note where it has one, otherwise picked by eye from six random events of the class.

\newcommand{\lossfig}[2]{\begin{figure}[p]\centering\includegraphics[width=\textwidth]{#1}\caption{#2}\label{fig:#1}\end{figure}}
\lossfig{loss_band}{Noise band: a whole-window baseline offset on a contiguous patch of strips (here the y view, right), seen in x and y together; the track itself is clean, the view is lost to the competing-cluster cut.}
\lossfig{loss_multi_m3picked}{Second particle: the reference particle is fitted, but a separate particle in the same view makes the view fail the competing-cluster cut.}
\lossfig{loss_multi_otherpicked}{Second particle: production fits the other, brighter particle.}
\lossfig{loss_multi_merged}{Second particle within the 12\,mm seed gap in one projection: the two tracks are merged and the fit lies between them.}
\lossfig{loss_spark_discharge}{``Spark'' (more than 50 strips above the seeding floor): a discharge, whole chips light up and ring.}
\lossfig{loss_spark_muonother}{``Spark'': a clean muon on the reference plus an unrelated deposit elsewhere in the plane.}
\lossfig{loss_spark_shower}{``Spark'': a shower developing at the crossing.}
\lossfig{loss_delta}{$\delta$-ray: a sideways deposit from the track pulls the straight-line fit.}
\lossfig{loss_satblob}{Saturated $\delta$ blob: a faint MIP line and a saturated deposit on the ADC rail (amber outline: samples $\ge$\,3550, censored; sky-blue outline: samples on the bottom rail after the clip, fitted as data). The fit tilts towards the deposit. Censoring at each channel's own rail ($4095-$pedestal$-$CM) and the bottom rail, or adding one point-like deposit to the forward model, moves agreement with the reference on 350 saturated views of Run~1, refitted on the two-track re-pass (reconstruction, calibration and reference as in the pass), from 71.4 to at most 72.9\,\% (best: both rails and the compact deposit; the multi-start fit alone gives 71.1); with the track pinned to the reference the same model fits worse in 93\,\% of the failures: the deposit is a short $\delta$ track over 5--10 strips, not a point (\texttt{rerun/satstudy/satblob.py}).}
\lossfig{loss_fit}{Fit off with no second particle-like cluster or saturation flagged (outlier class FIT).}
\lossfig{loss_headon}{Fit failed in one view: the X view is fitted on the reference, but in Y the track lies beside a patch of strips with a whole-window baseline offset near the passivated end and no Y fit is returned. This is the only such track in Run~1 on the re-pass. The head-on class of the table (a two-strip x cluster whose third strip is below the 10\,\% relative floor, never seeded) is recovered by the pass's head-on seeding, so no head-on loss is left to display.}
\lossfig{loss_postail}{Position tail: the fit agrees with the charge, the reference does not (about two thirds are tracks bent between the reference doublets).}
\lossfig{loss_edge}{Edge of the other coordinate.}
\lossfig{loss_multi_deposit}{A second deposit that is not particle-like.}

\analysis{A1}{decide whether B/E/C/D get a column; one paragraph per class; repeat the saturated-blob test on the re-pass or keep it marked pre-pass}
